\chapter{不安全代码与原始存储}
\label{ch:23}

不安全操作把某些证明义务交给程序作者。把它们限制在小函数内，写清存储状态、边界、对齐、所有权和清理条件，再向外提供普通类型接口。

\section{地址不是任意有效对象}

std::lang::address 查询地址表示，地址运算需要核对单位、偏移、有效分配和目标类型。数值相近不能证明两个地址属于同一对象；checked 的运行期检查也不能代替外部内存的所有前提。

\par
\medskip
\begin{lstlisting}[style=chali]
unsafe int64 addOffset(int64 address, int64 bytes) {
    return address + bytes;
}

void main() {
    unsafe {
        println("{}", addOffset(100, 8));
    }
}
\end{lstlisting}
\par
\medskip

\par
\medskip
\begin{lstlisting}[style=out]
108
\end{lstlisting}
\par
\medskip

\section{原始容量与初始化状态}

RawStorage<T> 分配容量但不构造全部元素。initAt 在未初始化槽位构造值，takeAt 移出并使槽位重新未初始化，dropAt 清理存活元素但保留容量。清理时不能把同一拥有值析构两次。

\par
\medskip
\begin{lstlisting}[style=chali]

void main() {
    RawStorage<String> slots = RawStorage<String>::allocate(4);
    unsafe {
        initAt<String>(slots.pointer(0), "hello");
        String result = takeAt<String>(slots.pointer(0));
        assert(!slots.isInitialized(0));
        println("{} / {}", result, slots.capacity());
    }
    slots.close();
}
\end{lstlisting}
\par
\medskip

\par
\medskip
\begin{lstlisting}[style=out]
hello / 4
\end{lstlisting}
\par
\medskip

\par
\medskip
\begin{lstlisting}[style=chali]

void main() {
    RawStorage<int64> slots = RawStorage<int64>::allocate(3);
    unsafe {
        BarePtr<int64> p = barePointer<int64>(slots.pointer(0));
        BarePtr<int64> q = p + 2;
        assert(q.address - p.address == 16);
        assert((2 + p).address == q.address);
        q = q - 1;
        println("{} / {}", q - p, q.address - p.address);
    }
    slots.close();
}
\end{lstlisting}
\par
\medskip

\par
\medskip
\begin{lstlisting}[style=out]
1 / 8
\end{lstlisting}
\par
\medskip

\section{复制字节与复制值}

memcpy 等原语只按其内存约定工作，不能普遍替代拥有型对象的复制操作。按字节复制整数和按字节复制 String 的拥有信息不是同一回事；后者可能破坏所有权、引用计数或清理状态。

\par
\medskip
\begin{lstlisting}[style=chali]

void main() {
    RawStorage<int64> source = RawStorage<int64>::allocate(1);
    RawStorage<uint8> target = RawStorage<uint8>::allocate(8);
    unsafe {
        BarePtr<int64> from = barePointer<int64>(source.pointer(0));
        BarePtr<uint8> to = barePointer<uint8>(target.pointer(0));
        memset<int64>(from, (uint8)42, 8);
        memcpy<uint8,int64>(to, from, 8);
        uint8 mut& first = dereference<uint8>(to);
        uint8 mut& last = dereference<uint8>(to + 7);
        println("{} / {}", (int64)*first, (int64)*last);
        assert(!source.isInitialized(0));
    }
    source.close();
    target.close();
}
\end{lstlisting}
\par
\medskip

\par
\medskip
\begin{lstlisting}[style=out]
42 / 42
\end{lstlisting}
\par
\medskip

\section{封装和验证不变量}

在安全入口验证数量和边界，在内部维护已初始化区间，在退出路径清理实际存活元素。rawSize 描述对应查询的原始存储大小，不能据此假定宿主 ABI、所有对齐或任意序列化格式。完整边界见 \href{https://github.com/rabbitc999/ChaLi/blob/master/doc/41-unsafe-abi.md}{不安全存储与 ABI}。

\section{习题}

\begin{enumerate}

\item takeAt 后还能直接 readAt 同一槽位吗？

\item dropAt 是否释放整个缓冲区？

\item 为何 memcpy 不能替代所有复制？

\item 封装原始存储时应记录哪些状态？

\end{enumerate}

\subsection*{参考解答与验收要点}

\begin{enumerate}

\item 不能，必须重新初始化。

\item 不会，清理该元素并保留容量。

\item 拥有值包含清理和所有权语义。

\item 容量、初始化区间、存活元素和借用有效期。

\end{enumerate}

下一章：\hyperref[ch:24]{第 24 章}。

\section{番外：油杓与葫芦}

改写自欧阳修《卖油翁》。

陈尧咨见老翁酌油，问道：“既然手熟，何必还要铜钱？”翁曰：“钱孔使我看清边界，手熟使我守住边界，二者并不相妨。”

有门客闻言，将铜钱收进袖里，又在葫芦上贴了“不安全”三字，自谓已尽得要领。油未入瓶，先湿了衣袖。公问其故，门客答：“既已声明不安全，自然可以不问油流向哪里。”

翁取杓而止之：“你可以亲手持杓，却不能把别人的葫芦当作自己的。倒进去的是油，不是瓶子的主人；抄下来的是瓶号，也不是第二只瓶子。”

于是门客把取油、持杓、换瓶各写一条规矩，替客人留下一个稳妥的把手。客人只须接过葫芦，如何守住钱孔，仍由持杓的人负责。公笑曰：“这回写在纸上的，终于不是免罪的告示。”

\subsection*{技术对照}

\texttt{unsafe} 表明实现者承担额外证明义务。裸地址和按位复制不会产生资源所有权，范围检查与清理协议仍各有职责。把底层操作封装成接口，是让调用者使用已经维护好的不变量；性能则要通过实际负载检验，不能由一个关键字推断。
